\begin{abstract}
Tree speculative decoding for recurrent-hybrid language models must preserve candidate ancestry and continuation state. We study GDN Tree-Scan, a Gated DeltaNet verifier whose accepted-prefix contract coordinates recurrent state, convolution history, attention storage, and pending correction tokens. Archived regression witnesses expose incomplete state publication and physical-layout-dependent numerical disagreement. The publication repair changes one token fixture from 15 failures in 15 checks to zero. Fresh same-input measurements compare sequential scan/replay with local triangular-solve and finite-Neumann reconstructions on eight pilot prefixes and 31 of 32 confirmation prefixes passing frozen provenance checks. Frozen numerical, padding, and timing criteria fail, preventing promotion of the compact routes; a retrospective audit also identifies omitted checks and inaccurate pilot rationale. After bounded state, greedy-path, and output checks, eighteen fresh boots on GB10 compare three instrumented Qwen3.6-27B-FP8 configurations. Against native five-step multi-token prediction, the tree averages 12.69 versus 13.67 tokens/s at batch one and 47.28 versus 55.38 at batch four; it exceeds the longer eleven-step control. These rates bind returned tokens to the same physical wall intervals. Engine seeds differ across arms, and only 14 of 144 same-block, across-arm prompt-stream comparisons are exactly equal. The results therefore describe the configurations as executed, with coarse three-block uncertainty. They provide scoped state-integrity witnesses and numerical/cost characterization, without establishing full-model sequential equivalence, quality preservation, or an isolated tree-algorithm speedup.
\end{abstract}
